\section{Extending localization to inspection outputs}\label{sec:r9-executor}
The service in Sections~\ref{sec:problem}--\ref{sec:theory} remains the numerical foundation. Its state, registration objective, source replacement, replay and computational policies are retained. The additional experiment exposes three consumer tasks: acquiring diagnostically usable synthetic pixels, associating them with the source version available at capture, and delivering a report within declared limits. These outputs test the operating contract beyond the localization proxy in Eq.~\eqref{eq:observation-event}; they do not replace the localization experiments or redefine their events.

The new catalogue contains Fixed2, Fixed4, Final, FlatCART, FlatRF, KG1, ContractRF and OldRAVC, with the unchanged fitted predictors, prices and common budget wrapper. One common solver prefix is actually executed per new input, and each policy/service wrapper consumes only the checkpoints it would have revealed. Its component timing is a paired controlled-prefix ledger; the complete-route timings in Section~\ref{sec:evaluation} remain separately measured online executions.

\subsection{Actual raster and diagnostic target processing}
A prescribed nominal route and prior CAD mesh plan up to 15 face/barycentric target addresses before stochastic map displacement and sensing. A synthetic gimbal aims at a prior-CAD patch. Targets remain the same across the paired ROI intervention. Raycasting against the perturbed geometric scene creates RGB pixels with procedural texture, fixed lighting, projected 25-mm-radius diagnostic landmarks, an exposure-blur surrogate and Gaussian intensity noise. A nearest-palette decoder estimates landmark pixel centers using only the image and known landmark colors. Ground truth is used for scene generation and offline correctness labels, not for decoding or online numerical policy choice. This is a diagnostic synthetic landmark task, not crack detection or a new real-camera experiment.

At a target with range $d$ and absolute ray/normal incidence $\gamma$, the defined geometric sampling density is $g=d/(f_{\rm px}\gamma)$. Fixed-target quality uses the maximum GSD and motion/exposure blur across the preplanned addresses, including targets outside the current ROI. Correctly acquired target coverage is scored from actual decoded pixels. The capture-time geotag projects the CAD targets using the causally installed localization state. Its pixel discrepancy and metric landmark displacement are separately scored. An absent positive source version fails the readiness requirement; a later update cannot retroactively count as capture-time ready.

Image capture occurs at an independent, before-sensing sampled tick $c\in\{0.1,0.2,0.3\}$~s. Report delay and image age are distinct:
\begin{equation}
t_{\rm report}=\max\{t_{\rm first\ positive},c+0.05+\alpha w_{\rm camera}\},\qquad
a_{\rm image}=t_{\rm report}-c.\label{eq:r9-report}
\end{equation}
The camera clock characterizes the executed \emph{virtual raster producer plus decoder}, not a physical camera's capture latency. Its excluded-pilot median is frozen separately for each ROI. The same declared load multiplier is used across methods. Native raster generation, decoding and localization components are reported separately; their sum is a mixed-boundary virtual workload, not embedded online CPU or worst-case timing. This narrow timing model is part of the certified simulation event and cannot be transferred directly to physical inspection.
